% year: תשע"ח
% type: מבחן
% semester: א
% moed: ג
% number: 3א
% points: 10
% categories: taylor, func-limits
% source: exams/מבחנים/תשעח/מועד מיוחד תשעח.pdf

\begin{question}
חשבו את הגבול:
\[ \lim_{x\to0}\left(\frac{1-\cos(x^3)}{x^2\ln(1+x^2)\sin^2x}\right) \]
\underline{רמז}: אפשר להשתמש בנוסחאות מקלורין.
\end{question}

\begin{hint}
$1-\cos u=\frac{u^2}{2}+o(u^2)$, $\ln(1+u)=u+o(u)$, $\sin x=x+o(x)$. מה סדר הגודל של המונה ושל המכנה?
\end{hint}

\begin{solution}
פיתוחי מקלורין: $\cos u=1-\frac{u^2}{2}+o(u^2)$ עם $u=x^3$ נותן
\[ 1-\cos(x^3)=\frac{x^6}{2}+o(x^6). \]
כמו כן $\ln(1+x^2)=x^2+o(x^2)$ ו-$\sin x=x+o(x)$, ולכן $\sin^2x=x^2+o(x^2)$. המכנה:
\[ x^2\ln(1+x^2)\sin^2x=x^2\left(x^2+o(x^2)\right)\left(x^2+o(x^2)\right)=x^6+o(x^6). \]
לכן
\[ \frac{1-\cos(x^3)}{x^2\ln(1+x^2)\sin^2x}=\frac{\frac12+\frac{o(x^6)}{x^6}}{1+\frac{o(x^6)}{x^6}}\xrightarrow[x\to0]{}\frac12. \]
(באופן שקול: $\frac{1-\cos(x^3)}{x^6}\to\frac12$, $\frac{\ln(1+x^2)}{x^2}\to1$, $\frac{\sin^2x}{x^2}\to1$, ומפרקים את הביטוי למכפלה.)
\[ \boxed{\lim_{x\to0}\frac{1-\cos(x^3)}{x^2\ln(1+x^2)\sin^2x}=\frac12} \]
\end{solution}
